\section{Durability}
\label{sec:durability}

Process monitoring is high-volume and recent-state oriented. Data-individual
metadata capture has a stronger requirement: it must not be silently lost. This
section describes the append-only segment substrate, the process-event journal,
and the metadata outbox that together provide durable, replayable local capture.

\subsection{Segment Format}

Segments are a dependency-free durability substrate for raw events and metadata
outbox entries. The format is intentionally small: an eight-byte magic header
followed by length-prefixed, checksummed records.

\begin{lstlisting}[numbers=none]
magic[8]                  // the bytes "VANNAK01"
repeated:
  len: u32 little endian
  checksum: u64 little endian
  payload: [u8; len]
\end{lstlisting}

The per-record checksum is a stable non-cryptographic FNV-1a derivative over the
payload, for corruption detection rather than tamper resistance. A rolling
segment checksum combines record checksums with
\texttt{current.rotate\_left(7) \^{} next}.

\subsection{SegmentWriter and SegmentReader}

\texttt{SegmentWriter::create} opens a new file (failing if it already exists),
writes the magic header, and tracks \texttt{record\_count}, \texttt{byte\_len},
and the rolling checksum. \texttt{append\_record} returns the
\texttt{RecordOffset} at which the record begins; records larger than
\texttt{u32::MAX} are rejected with \texttt{SegmentError::RecordTooLarge}.

\begin{lstlisting}
pub fn append_record(&mut self, payload: &[u8]) -> Result<RecordOffset, SegmentError> {
    if payload.len() > MAX_RECORD_LEN {
        return Err(SegmentError::RecordTooLarge { len: payload.len() });
    }
    let offset = self.byte_len;
    let record_checksum = checksum(payload);
    let len = payload.len() as u32;
    self.file.write_all(&len.to_le_bytes())?;
    self.file.write_all(&record_checksum.to_le_bytes())?;
    self.file.write_all(payload)?;
    self.record_count += 1;
    self.byte_len += 4 + 8 + payload.len() as u64;
    self.checksum = combine_checksum(self.checksum, record_checksum);
    Ok(RecordOffset(offset))
}
\end{lstlisting}

\texttt{SegmentReader::open} validates the magic (returning
\texttt{SegmentError::InvalidMagic} on mismatch). \texttt{read\_next} returns
\texttt{Ok(None)} at end of file and verifies each record's checksum, returning
\texttt{SegmentError::ChecksumMismatch} when it differs.

\texttt{SegmentWriter::open\_append} validates an existing segment by scanning
all records, rejects corrupt or trailing partial records, and reopens the file
for append while preserving \texttt{record\_count}, \texttt{byte\_len}, and the
rolling checksum. This is the primitive used by both process-event and metadata
outbox recovery.

\subsection{SegmentManifest}

A sealed or in-progress segment is described by an owned
\texttt{SegmentManifest}:

\begin{lstlisting}
pub struct SegmentManifest {
    pub segment_id: SegmentId,
    pub node_id: NodeId,
    pub path: PathBuf,
    pub record_count: u64,
    pub byte_len: u64,
    pub checksum: u64,
}
\end{lstlisting}

\texttt{SegmentWriter::seal} flushes, syncs, and returns the final manifest. The
manifest is the unit that is later recorded into the cluster control state as a
sealed segment.

\subsection{ProcessEventJournal}

\texttt{ProcessEventJournal} is the typed local history for accepted
\texttt{PipelineEvent} values. It uses the segment format above with a compact
binary payload codec for the process-event envelope, including Durga lifecycle
kind/status, tenant and environment identity, process instance identity, optional
activity, correlation, error, metadata references, causal parent, and payload
fields.

\texttt{ProcessEventJournal::append\_durable} validates the event, appends the
encoded payload, and syncs the segment before returning the record offset. The
service-level durable ingest path,
\texttt{VannakService::ingest\_process\_event\_durable}, performs this journal
write before reducing the event into the hot index.

\texttt{replay\_process\_event\_segment} reads a segment back into owned
\texttt{JournaledPipelineEvent} values and reports a
\texttt{ProcessEventReplaySummary}. \texttt{ProcessEventJournal::recover}
combines replay with validated append reopening, and
\texttt{VannakService::recover\_process\_events\_into\_hot\_index} can rebuild a
single-node hot index from the local process-event segment.

\subsection{MetadataOutbox}

\texttt{MetadataOutbox} is the in-memory delivery queue. Entries are keyed by
\texttt{IdempotencyKey}, with a separate \texttt{pending} queue for ordering.
\texttt{enqueue} returns \texttt{OutboxEnqueueResult::Enqueued} for a new key or
\texttt{Duplicate} for a known one. Each \texttt{MetadataOutboxEntry} carries its
payload, status, optional record offset, retry count, and last error.

\subsubsection{Pending, Acknowledged, Failed}

\texttt{OutboxStatus} has three states. \texttt{acknowledge} marks an entry
delivered and removes it from the pending queue. \texttt{fail} marks it failed,
increments \texttt{retry\_count}, records the error, and removes it from
pending. \texttt{retry\_failed} returns a failed entry to pending.

\subsection{Delivery}

\texttt{deliver\_next\_pending} takes the next pending entry, calls the
\texttt{IptoWriter}, and either acknowledges it or records the failure (carrying
the writer's retryable flag).

\begin{lstlisting}
pub enum MetadataOutboxDeliveryResult {
    NoPending,
    Acknowledged { idempotency_key: IdempotencyKey },
    Failed { idempotency_key: IdempotencyKey, retryable: bool, message: String },
}
\end{lstlisting}

\texttt{drain\_pending\_outbox} delivers up to \texttt{max\_attempts} entries and
returns a \texttt{MetadataOutboxDrainSummary} with \texttt{attempted},
\texttt{acknowledged}, \texttt{failed}, and \texttt{stopped\_after\_failure}.
Draining stops early after a retryable failure so transient backend outages do
not burn the attempt budget.

\subsection{Segment-Backed Outbox}

\texttt{SegmentBackedMetadataOutbox} couples an outbox with a
\texttt{SegmentWriter}. The key method, \texttt{enqueue\_durable}, appends and
syncs the encoded payload to the segment \emph{before} making the entry visible
as pending in memory. This guarantees that capture is durable before it is ever
acknowledged.

\begin{lstlisting}
pub fn enqueue_durable(&mut self, payload: IptoWritePayload)
    -> Result<DurableOutboxEnqueueResult, MetadataOutboxStorageError>
{
    if self.outbox.entries.contains_key(&payload.idempotency_key) {
        return Ok(DurableOutboxEnqueueResult::Duplicate);
    }
    let offset = self.writer.append_record(&payload.encode())?;
    self.writer.sync()?;
    self.outbox.enqueue_with_offset(payload, Some(offset));
    Ok(DurableOutboxEnqueueResult::Enqueued { offset })
}
\end{lstlisting}

\begin{vnote}
Durable capture is never acknowledged before the encoded payload is persisted to
a segment. This honors the principle that data-individual provenance must not be
treated as best-effort monitoring.
\end{vnote}

\subsection{Snapshots}

\texttt{MetadataOutbox::snapshot} returns a \texttt{MetadataOutboxSnapshot} with
\texttt{total}, \texttt{pending}, \texttt{acknowledged}, \texttt{failed}, and
\texttt{queued\_pending} counts. The segment-backed wrapper returns a
\texttt{SegmentBackedMetadataOutboxSnapshot} that pairs the outbox snapshot with
the current \texttt{SegmentManifest}.

\subsection{Segment Replay}

After a restart, pending delivery state is rebuilt from the segment.
\texttt{replay\_metadata\_outbox\_segment} decodes every record into a fresh
outbox. \texttt{replay\_metadata\_outbox\_segment\_after} skips records at or
before a checkpoint offset and returns a \texttt{MetadataOutboxReplay} containing
the rebuilt outbox and a \texttt{MetadataOutboxReplaySummary}.

\begin{lstlisting}
pub struct MetadataOutboxReplaySummary {
    pub checkpoint_offset: Option<RecordOffset>,
    pub scanned_records: usize,
    pub skipped_records: usize,
    pub replayed_records: usize,
}
\end{lstlisting}

\subsection{Checkpoints}

\texttt{SegmentBackedMetadataOutbox::acknowledged\_checkpoint} produces a
\texttt{MetadataOutboxCheckpoint} from the highest acknowledged record offset for
a given target. This checkpoint can then be proposed through the cluster control
plane so a recovering node can replay only the records after it.

\subsection{Outbox Checkpoint Data}

\begin{lstlisting}
pub struct MetadataOutboxCheckpoint {
    pub data_individual_shard_id: DataIndividualShardId,
    pub target: IptoInstanceId,
    pub segment_id: SegmentId,
    pub last_acknowledged_offset: RecordOffset,
    pub mapping_version: String,
    pub epoch: CheckpointEpoch,
}
\end{lstlisting}

The checkpoint records which shard and target it covers, the segment and offset
up to which delivery is durable, the mapping version in force, and a monotonic
\texttt{CheckpointEpoch} used to reject stale records.

\endinput
